% Source: 62_Chapter_48_Windows_to_the_Christian_World_Places_a.txt
\bookchapter[Windows to the Christian World]{CHAPTER 48}{WINDOWS TO THE CHRISTIAN WORLD}{Places and Persons of Faith}

\leadin{W}{E ARE HUMBLED TO} remember that we describe the Global Church in Western terms. As Randy
Richards has written, “generalizations are always wrong and usually helpful.” Perhaps some texture
can be added to our description and exposition of the Global South by encountering the stories,
however briefly, of a few places and the people who together comprise the great people of God.

\section{CHINA}

We will consider first and at greatest length the amazing story of faith in China. China had known
great hostility over religious turmoil. One religious uprising led to twenty-five million deaths!
Christian missionary efforts had been received with suspicion and resulted in an expulsion of
missionaries in the ninth and thirteenth centuries. The Jesuits had more enduring but modest success
in the 1600s. Still other Christians complained that the Jesuits had compromised Christian teachings
in order to make them accessible to the Chinese.

Western nations made China the focus of their missionary work in the late 1800s and early 1900s.
After World War II, a civil war broke out between the Nationalist Party led by Chiang Kai-shek and
the Communist Party led by Mao Zedong. Eventually the Nationalist Party was pushed off the mainland,
settling in Taiwan. The Communists formed The People’s Republic of China and expelled the Christian
missionaries in 1949, despite their long tenure in the country. Many wondered whether Christians
would survive, given China’s past violence over religion, the suspicion that Christianity was a tool
of Western imperialism, the weak, nominal state of some converts, and the Communists’ hostility to
religion.

At first there was an effort to control Christianity and other religions. Mao Zedong believed that
the faith would die a natural death when the oppression and manipulation of religion was restrained.
Governing bodies were created to ensure patriotic compliance in matters of religion. The Protestant
Three-Self Patriotic movement (1955) and the Catholic Patriotic Association (1957) became powerful
bodies that censured and coerced compliance in the name of loyalty to the nation. Chinese Christians
were badly divided. Some joined and endured difficulties to ensure that a Christian voice would
survive in China. Others protested that such cooperation was a compromise of the faith.

The Protestant organization took its name from an old missionary strategy that called for churches
to be more self-reliant: “self-governing; self-supporting, and self-propagating.” The co-opted terms
had additional meaning for Communists, who saw Western influence as necessarily oppressive. The
Protestants who cooperated were viewed as traitors. They were called the “3 self-church,” but more
commonly now “the registered Church.” The others risked life and limb to maintain independence. They
were called “the underground church,” “the house church,” and more commonly today, “the unregistered
church.” The Chinese government’s policy included brutal persecution, but it also believed that
isolating Chinese Christianity from Western influence would lead to its collapse. A later Chinese
policy attempted to bring religion to a brutal end. Chairman Mao himself instigated the Great
Proletarian Cultural Revolution (1966–1976). Massive brutality, re-education, and cult-like
veneration of Mao Zedong were the norm. Church property was confiscated, and even those Christians
who had cooperated with the registered churches were persecuted.

After Mao Zedong’s death, more moderate voices gained influence in China. Deng Xiaoping embodied
this more open approach. President Jimmy Carter normalized relations with China in January of 1979.
Visiting the United States, Deng Xiaoping thanked Carter for benefiting the people of China, and
asked if he could extend a kindness toward Carter. Carter told him that as a young boy he had given
five cents a week to build hospitals and schools in China and had viewed missionaries as heroes of
his faith. Carter observed that China forbade missionaries, Bibles, and freedom of worship, and
requested that all three be allowed. Deng Xiaoping asked for time to think over the requests. He
answered Carter the next morning. He said that China would not allow missionaries but would
authorize the distribution of Bibles and would pass a law allowing free worship. Permitting the
printing and distribution of Bibles came more quickly, and in 1982 a law began the process of
allowing Christians greater freedom to worship. The actual freedom to worship varies greatly from
region to region, but more and more the unregistered (underground) churches are holding open
meetings and the registered churches are enjoying greater liberty.

The lines of hostility are diminishing, with cooperation between the legal registered church and the
unregistered churches becoming more typical. Another important voice and influence has emerged in
addition to the registered and unregistered churches. China’s economic success has brought
entrepreneurs from around the globe. Internationals with passports are given the legal right to free
worship and assembly. These churches composed of Christians from across the world form an important
spiritual influence and liaison to the other churches in China.

The number of converts in China is mind-bending. A conservative estimate for 2000 is 89 million
believers, a full 7.1 percent of the Chinese population.

\section{CARDINAL KUNG: FAITHFUL RESOLVE}

Ignatius Cardinal Kung, Kung Pin-Mei (also Gung Pin Mei) was ordained in 1930 and in 1950 became the
first native-born Chinese bishop of the important port city of Shanghai. He was imprisoned in 1955
for his unwillingness to renounce the pope and recognize the legitimacy of Chinese efforts to
control the Catholic Church. The authorities paraded the bishop before a huge audience at the Dog
Racing Stadium to denounce the pope. Instead, he defiantly shouted, “Long live Christ the King, long
live the pope.” He was imprisoned for the next five years before being convicted of treason in 1960.
For the next twenty-five, years he was isolated and kept alone except for times of physical labor.
He was not given a Bible or a rosary. He was not allowed correspondence or contact with his family.

After the death of Chairman Mao, China was moving slowly toward a more open economy and, to a lesser
and slower extent, a more open culture. Deng Xiaoping arranged a banquet for the cardinal and bishop
of Manila, Jaime Sin, in 1984. This Filipino churchman had been instrumental in the popular uprising
that unseated Ferdinand Marcos. As a gesture to this openness, the old bishop, who had been in
prison for almost thirty years, was invited to the banquet. The carefully marshaled circumstances
did not permit the two churchmen to share a personal or private moment. Still, the old bishop seized
the opportunity to share a song. Defiantly he sang a song in Latin based on Matthew 16:18: “Upon the
rock I will build my church.” His resolve had survived thirty years of labor and isolation. The next
year, the Chinese would release Kung for “health reasons.” Not even Cardinal Sin knew that evening
that Pope John Paul II had already named Kung as a cardinal but had done so in pectore (literally in
the heart, a secret designation when publicity might be inappropriate or even dangerous). Kung’s
life runs a remarkable course. He attempted to work with Chinese authorities but later defied them.
Thirty years of isolation and labor did not silence the resolve of this believer. Many worried that
the Jesuits had compromised Christianity; others wondered why any genuine believer would ever work
with the registered church. Yet time has given evidence of the Spirit’s soul-building to fashion
strong faith from humble beginnings.

\section{KOREA}

Korea’s encounter with Christianity is atypical. With some qualification, Koreans actually brought
Christianity to Korea. In 1784, Yi Singhun was baptized while on a diplomatic mission to China. He
returned with literature by the famous Jesuit missionary Mateo Ricci. Despite not having a priest,
the Korean Christians baptized each other and held mass. In 1794, a priest from China, James Chou,
was secretly sent to Korea where he took up ministering to almost four thousand practicing
Christians. Persecution followed these Catholics who were feared to have alliances with foreign
power; waves of martyrdom, including Chou himself, continued for nearly half a century.

In 1873, King Ko-jong displayed a greater openness toward western influence. In this new era,
Catholics grew in number. A Scottish Presbyterian missionary, John Ross, translated the New
Testament for the Koreans while serving in China. By the 1880s, numerous Protestant missionaries
began work in Korea.

The Protestant work reached a young man named Sun Chu Kil. His family moved to Pyongyang in 1885.
That same year, several American Presbyterian missionaries arrived. These missionaries supported
Bible translation and employed the strategy of John Nevius, an American missionary to China. Nevius
supported strong, self-reliant churches. The missionaries’ arrival corresponded with the Korean’s
disillusionment over its political circumstance. Korea was forced to exist in the shadow of China,
Russia, and Japan. Kil had struggled with his Taoism and begun to explore Christianity. One of the
American missionaries, Samuel Moffett, invited Kil to read Pilgrim’s Progress. During a season of
spiritual struggle, Kil spent a night weeping and praying and was converted.

Kil took an active part in church leadership and evangelism among the working class. In 1903, he
entered Pyongyang Presbyterian Seminary (supported by Moffett). Reports of the Welsh revival
provoked Kil to hold morning prayer meetings where they petitioned for a Korean revival. Near the
end of a ten-day training series for new converts, a missionary publically confessed “stubbornness
and pride.” Soon many others began to publically confess. Kil himself acknowledged his
inappropriately handling a will. The meeting lasted from eight in the evening until two in the
morning. The Koreans prayed simultaneously; each believer prayed aloud at the same time. Some say
the many spoken voices make a symphony pleasing to God. Numerous Westerners have observed this as a
distinctive experience of the Korean spirituality (though it is done elsewhere).

In 1910, the Japanese annexed Korea and ruled over it until the end of World War II. They were
brutal toward the churches. Many Christians led in advocating independence from Japan. Kil was a
powerful voice in the movement, though he maintained his embrace of Christian nonviolence. For once,
Christians earned a reputation for standing against the forces of colonization. In time, Kil was
appointed to an influential pastorate: Central Presbyterian Church. His leadership in independence
eventually cost him his life and the life of his son. He turned himself in to authorities after his
peers in the independence movement were arrested; after his release, he died of complications
resulting from his mistreatment. Mark Noll reports that Kil preached over seventeen thousand sermons
and shared in starting sixty different churches

After World War II, Korea was divided between the Soviet-supported North and the South, supported by
the United States. The fighting ceased in 1953, but the brutality towards Christians remained in
place. The Communists struck out against the church. At that time, two-thirds of Korean Christians
were in the North. It is very difficult to assess the state of Christianity in North Korea because
of the extremely secretive posture of the government. In South Korea, several Christian groups have
flourished. Presbyterianism, Pentecostalism, and, more recently, Catholicism have all grown
remarkably. Consequences of the Korean Revival were a catalyst to explosive growth. In 1914, one in
one hundred Koreans was Protestant. By 2010, one in three Koreans was Christian.

Korea is home to the world’s largest churches. The Yoido Full Gospel Church has well over a million
members. Its cell structure serves as model and perhaps inspiration to churches around the world.
South Korea follows only the United States in the formation of seminaries and sending missionaries
around the globe. These missionaries take with them the policy of building self-reliant churches,
but also share a Korean spirituality that is different in form and character. This distinctive
spirituality is rooted in prayer. Visitors report that Korean simultaneous prayer stands out. The
centrality of prayer is seen in the Prayer Mountain. Koreans book time at retreat centers for the
purpose of prayer and fasting. Pastor David Yonggi Cho of the Yoido church credits the growth of the
church, not to its cell structure or its amazing efficiency, but to its prayer.

\section{EAST AFRICAN REVIVAL}

The stories of two men and the friendship they forged is one piece of the great revival to sweep
across East Africa. Both of these men suffered discouragement. Simeon Nsibambi was born in 1897. He
was converted in 1922. Numerous opportunities helped create the hope and expectation that he would
be selected to study abroad, but it was not to be. God addressed Simeon’s severe disappointment in a
vision. In the vision, God asked him to compare the value of the scholarship with the value of the
gospel message and the forgiveness that had been granted him.

Dr. Joseph Church was converted as a student at Cambridge in August, 1920. Cambridge had known
evangelical stirrings. American evangelist D. L. Moody had preached at Cambridge almost four decades
before Church’s experience of conversion. In the wake of Moody’s meeting, numerous Cambridge
students were called to serve in missions, including the Cambridge Seven who targeted China. One
such student, George Lawrence Pilkington, went to Uganda where he witnessed a revival in 1893, a
precursor to the revival to follow. Dr. Church applied his skills as a doctor in Africa where he
also distinguished himself as an able missionary.

By 1929, Church was undergoing a serious physical and spiritual discouragement. He longed for
companionship to address his emptiness when he encountered Simeon Nsibambi. The two shared their
longings for revival for themselves and Africa. They prayed and studied the Bible together for two
days. The experience and unity of purpose was life changing. Noll cites Church’s reflections upon
the meeting “as the time that God in his sovereign grace met with me and brought me to the end of
myself and thought fit to give me a share in the power of Pentecost.”

The fellowship that included these two would grow to be extended again and again. In September,
1935, another group met for spiritual renewal; their thinking was shaped by the Keswick theme that
complete surrender to Christ preceded Spiritual breakthroughs. They planned and initiated a series
of meetings that may be identified as the revival’s starting point. The group included Joseph Church
(the only Anglo) and Simeon Nsibambi. William Nagenda soon joined the circle and would become more
influential than any other group member in the revival’s spread. Two important words reveal the
revivalist nature of this outpouring: abaka means on fire; balokole means the saved.

Some credit this core leadership not only for preaching but for the organizational wisdom to send
groups out in teams that involved lay persons. The revival’s geographic sweep moved from North
Uganda to include Kenya, Rwanda, Burundi, and Tanzania (then Tanganyika). It touched many
internationally beyond the region. The numbers defy all imagination. The percentages of Christian
adherents near the century’s end are as follows: Uganda 89 percent, Kenya 79 percent, Rwanda 81
percent, Burundi 90 percent, and Tanzania 51 percent.

William Nagenda and Joe Church made perhaps the most effective team. Nagenda’s strong family was a
vindication to his powerful preaching. Nagenda also was involved in one of the important tests faced
during the revival. In 1941 twenty-nine students were expelled from Bishop Tucker Theological
College. The school’s leadership were alarmed at the revival zeal of these balokole, or saved ones.
Thankfully the (all Anglo) school authorities were able to restore almost every student, including
William Nagenda. Students were able to express their revival fervor and remain loyal to the
denomination; this flexibility proved very important for the revival movement. The movement,
however, crossed denominational lines at numerous points. Brethren, Seventh Day Adventists,
Mennonites, Baptists, and Methodists all report participation in the revival.

Another celebrated voice of the movement is Festo Kivengere. Kivengere was born to a ruling family
and worked as a shepherd as a young boy. The young shepherd read children’s stories about Jesus as
he worked. He was converted during an evangelistic meeting. He studied theology, pastored a church,
and was eventually ordained as bishop in the Anglican Church. He befriended Billy Graham and served
as his translator as well as appearing at Graham’s crusades. Kivengere’s courage encountering Idi
Amin was widely admired. Amin murdered his archbishop, Janani Luwum, and Kivengere remained in exile
until Amin’s ouster. During this time, he became an internationally recognized voice for social
justice and for all of Africa. His theology was firmly rooted in the cross. Kivengere believed that
Christ called him to love Amin despite his atrocities.

In China, Korea, and East Africa we observe several common elements. Each has an array of its own
rich Christian cultures while having indebtedness to missionary efforts. Each has seen remarkable
numerical growth. Each has produced believers who give to Christ in their bearing suffering and
extending grace.

\readinghead{Suggestions for Further Reading}
\begin{readinglist}
\entry *Anderson, Allen. To the Ends of the Earth: Pentecostalism and the Transformation of World
  Christianity. New York: Oxford University Press, 2013.
\entry *Hill, Jonathan. Zondervan Handbook to the History of Christianity. Grand Rapids: Zondervan,
  2006.
\entry *Jacobsen, Douglas. The World’s Christians: Who they Are, Where They Are, and How They Got
  There. Oxford, UK: Wiley-Blackwell, 2011.
\entry *Osborn, H. H. Pioneers in the East African Revival. Winchester, UK: Apologia, 2000.
\entry *Noll, Mark A. and Carolyn Nystrom. Clouds of Witnesses: Christian Voices from Africa and
  Asia. Downer Grove, IL: InterVarsity Press, 2011.
\entry *Spickard, Paul R. and Kevin M. Cragg, A Global History of Christians: How Everyday Believers
  Experience Their World. Grand Rapids: Baker, 2004.
\end{readinglist}
